\subsection{\textit{Sprint} 1}

Para esta \textit{sprint}, as atividades previstas para a equipe eram fazer uma tela
de \textit{login} e uma tela \textit{home}. Quanto às atividades previstas para meu \textit{squad}, estava
prevista a implementação do \textit{front-end} da tela de \textit{login}, utilizando React
\cite{react-docs} e TypeScript \cite{typescript-docs},
utilizando componentes genéricos criados pelo grupo, como um botão, por
exemplo, e em alguns casos até criando componentes genéricos, como \textit{inputs}.

A equipe inteira conseguiu completar as \textit{tasks} e a parte do \textit{back-end}
conseguiu avançar até mais do que o previsto. Quanto ao meu \textit{squad}, conseguimos
implementar a tela de \textit{login} quase igual ao protótipo do Figma \cite{figma-help} e criar um novo
componente genérico de \textit{input}.

Por incrível que pareça não foram encontrados tantos problemas, creio que o
fato de a tela desenvolvida ser simples facilitou o desenvolvimento sem muitas
barreiras, apenas com algumas correções para garantir que a página fosse
responsiva e o mais genérica possível.

Quanto às lições, aprendi que a maioria dos itens utilizados na página devem
ser transformados em componentes com o intuito de facilitar o trabalho dos outros
colegas e até meu trabalho no futuro.

Prosseguirei, então, com minha equipe, que está planejando fazer cerca de 8
telas no \textit{front-end} e todo o \textit{back-end} por trás das que ainda não foram feitas (parte
destas 8 telas já tem o \textit{back-end} pronto). Quanto ao meu planejamento, estou
esperando meu \textit{squad} receber uma \textit{task}; no entanto, creio que serão cerca de 3 a 4 telas para meu \textit{squad} desenvolver.
